\documentclass[12pt]{article}
\usepackage[utf8]{inputenc}
\usepackage[T1]{fontenc}
\usepackage{hyperref}
\title{Yangum Dey Language Resources for K-12 Curriculum Development}
\author{Compiled by Claude Sonnet 5.5 (claude-sonnet-5-5) for Andr\'as Kornai}
\date{}
\begin{document}
\maketitle

\section{Introduction}
Yangum Dey (ISO 639-3 yde) is a Torricelli language of Sandaun Province, Papua New Guinea, one of three Yangum varieties (the Gel variety is reported to be nearly extinct). Speaker figures conflict. ASJP lists about 180 speakers, Joshua Project gives a population of 400 for the people group, and Wikipedia gives 960 for the wider Yangum cluster (2000 census). Very little exists for this language, and the search found no educational materials.

\section{Primary Educational Resources}
None were found.

\section{Linguistic and Reference Materials}
\begin{itemize}
\item Wikipedia, ``Yangum language''. \url{https://en.wikipedia.org/wiki/Yangum_language}
\item Glottolog entry for ISO code yde. \url{https://glottolog.org/resource/languoid/iso/yde}
\item ASJP wordlist entry, Yangum Dey. \url{https://asjp.clld.org/languages/YANGUM_DEY}
\end{itemize}

\section{Religious and Community Resources}
\begin{itemize}
\item Joshua Project, Yangum Dey people group. It is largely generic and unsourced, so its cultural details should be treated with caution. \url{https://joshuaproject.net/people_groups/18977}
\end{itemize}

\section{Digital and Audio Resources}
None were found.

\section{Curriculum Development Context}
There is no known schooling, orthography or published material in the language. Any curriculum would have to begin with community-led documentation and orthography development, with a speaker-led approach and in partnership with local education authorities.
\end{document}